\documentclass[10pt,a4paper]{article}
\usepackage[UTF8]{ctex}
\usepackage{amsmath,amssymb,geometry,xcolor,listings,tcolorbox,multicol,enumitem}
\geometry{margin=1.2cm,top=1cm,bottom=0.8cm}
\setlength{\parskip}{0pt}\setlength{\parsep}{0pt}\setlength{\itemsep}{0pt}
\setlist{nosep,leftmargin=*}

\lstset{language=C,basicstyle=\ttfamily\scriptsize,keywordstyle=\color{blue}\bfseries,
  numbers=none,frame=single,breaklines=true,showspaces=false,tabsize=2,
  aboveskip=2pt,belowskip=2pt}

\newtcolorbox{keybox}[1][]{colback=yellow!3!white,colframe=yellow!70!black,
  title=\textbf{考点},fonttitle=\bfseries\small,size=small,top=2pt,bottom=2pt,#1}
\newtcolorbox{bugbox}[1][]{colback=red!3!white,colframe=red!60!black,
  title=\textbf{易错点},fonttitle=\bfseries\small,size=small,top=2pt,bottom=2pt,#1}

\title{\textbf{线性结构：链表 · 栈 · 队列}}
\author{}\date{}

\begin{document}
\maketitle\vspace{-8pt}

\section*{链表 (Linked List)}

\textbf{核心思想：} 不需要连续内存。每个节点存数据 + 下一节点的指针。\\
\textbf{对比数组：} 数组 $O(1)$ 查找但 $O(N)$ 插入/删除；链表 $O(N)$ 查找但 $O(1)$ 插入/删除。

\begin{multicols}{2}

\subsection*{节点结构}

\begin{lstlisting}
typedef struct node {
    int data;
    struct node *next;
} Node;
\end{lstlisting}

\subsection*{创建节点}

\begin{lstlisting}
Node* createNode(int x) {
    Node* p = (Node*)malloc(sizeof(Node));
    p->data = x;
    p->next = NULL;    // 必须初始化为NULL！
    return p;
}
\end{lstlisting}

\subsection*{建链：串3个节点}

\begin{lstlisting}
Node* head = createNode(10);
head->next = createNode(20);
head->next->next = createNode(30);
// 10 -> 20 -> 30 -> NULL
\end{lstlisting}

\subsection*{遍历 (Traversal)}

\begin{lstlisting}
void printList(Node* head) {
    Node* p = head;           // 用临时指针
    while (p != NULL) {
        printf("%d ", p->data);
        p = p->next;
    }
}
\end{lstlisting}

\begin{bugbox}
\texttt{Node* p;} 只管声明不管分配 → 野指针。\\
遍历用 head 直接移动 → head 丢了再也找不到链表头。
\end{bugbox}

\subsection*{插入 (Insert)}

\textbf{在 node 后面插入：} 注意顺序！先连新节点的 next，再改 node 的 next。

\begin{lstlisting}
void insertAfter(Node* node, int x) {
    Node* p = createNode(x);     // 先建节点
    p->next = node->next;        // ① 新连后
    node->next = p;              // ② 前连新
    // 顺序反了链表会断！
}
\end{lstlisting}

\textbf{头部插入：} 需要改 head 本身，传 \texttt{Node**}

\begin{lstlisting}
void insertHead(Node** head, int x) {
    Node* p = createNode(x);
    p->next = *head;
    *head = p;           // 改main里的head
}
\end{lstlisting}

\begin{bugbox}
插中间时先 \texttt{node->next=p} 后 \texttt{p->next=node->next} → p 指向自己形成环！\\
头部插入用 \texttt{Node*} 传参 → 改不了 main 的 head。
\end{bugbox}

\subsection*{删除 (Delete)}

删第一个值为 x 的节点。分两情况：删头 vs 删中间。\textbf{一定要 free！}

\begin{lstlisting}
void deleteNode(Node** head, int x) {
    // Case 1: 删头节点
    if ((*head)->data == x) {     // 注意括号！
        Node* tmp = *head;
        *head = (*head)->next;
        free(tmp); return;
    }
    // Case 2: 删中间节点
    Node* p = *head;
    while (p->next && p->next->data != x)
        p = p->next;              // 必须判p->next非NULL
    if (p->next == NULL) return;  // 没找到
    Node* tmp = p->next;
    p->next = p->next->next;      // 绕过被删节点
    free(tmp);
}
\end{lstlisting}

\begin{bugbox}
\texttt{*head->data} 优先级错 → \texttt{*(head->data)}，应为 \texttt{(*head)->data}。\\
NULL 检查放在 \texttt{p->next->data} 后面 → 先崩了才判断。\\
忘了 free。\\
\texttt{x} 不存在时 \texttt{p->next->next} 段错误。
\end{bugbox}

\subsection*{双向循环链表}

\begin{lstlisting}
typedef struct node {
    struct node *llink;  // 指向前一个
    int item;
    struct node *rlink;  // 指向后一个
} Node;
\end{lstlisting}

性质: $\text{ptr} = \text{ptr}\to\text{llink}\to\text{rlink} = \text{ptr}\to\text{rlink}\to\text{llink}$

\subsection*{多项式 ADT}

链表表示：$\langle e_i, a_i \rangle$ 按指数降序。\\
系数数组浪费空间（稀疏多项式），链表节省空间。

\subsection*{游标实现 (无指针)}

全局数组 \texttt{CursorSpace[]}，索引 0 是自由链表头。\\
\texttt{malloc:} \texttt{p=CursorSpace[0].Next; CursorSpace[0].Next=CursorSpace[p].Next;}\\
\texttt{free:} \texttt{CursorSpace[p].Next=CursorSpace[0].Next; CursorSpace[0].Next=p;}

\columnbreak

\section*{栈 (Stack) — LIFO}

\textbf{后进先出。} 操作：Push(入栈)、Pop(出栈)、Top(查看栈顶)。

\subsection*{数组实现}

\begin{lstlisting}
#define MAX 100
typedef struct { int data[MAX]; int top; } Stack;
void init(Stack* s) { s->top = -1; }
void push(Stack* s, int x) {
    if(s->top>=MAX-1)return;
    s->data[++(s->top)] = x;       // 先+后放
}
int pop(Stack* s) {
    if(s->top==-1) return -1;
    return s->data[(s->top)--];     // 先取后-
}
int isEmpty(Stack* s) { return s->top==-1; }
\end{lstlisting}

\begin{bugbox}
push 忘判满，pop 忘判空。\\
\texttt{data[s->top++] = x} 后自增错 → \texttt{data[-1]}。
\end{bugbox}

\subsection*{栈应用一：后缀表达式求值}

遇到数字 push，运算符 pop 两个、算完 push 回去。\\
例：\texttt{6 2 / 3 - 4 2 * +} = 8。

\begin{tabular}{|c|c|c|}
\hline
\textbf{扫描} & \textbf{操作} & \textbf{栈} \\
\hline
6 & push 6 & [6] \\
2 & push 2 & [6,2] \\
/ & pop2,pop6→6/2=3,push & [3] \\
3 & push 3 & [3,3] \\
- & pop3,pop3→3-3=0,push & [0] \\
4 & push 4 & [0,4] \\
2 & push 2 & [0,4,2] \\
* & pop2,pop4→4*2=8,push & [0,8] \\
+ & pop8,pop0→0+8=8,push & [8] \\
\hline
\end{tabular}

\subsection*{栈应用二：中缀转后缀}

\begin{center}
\begin{tabular}{|c|c|c|}
\hline
\textbf{运算符} & \textbf{栈内 ISP} & \textbf{栈外 ICP} \\
\hline
\tt ( & 0 & 6 \\
\tt + - & 3 & 2 \\
\tt * / & 5 & 4 \\
\tt ) & — & 1 \\
\hline
\end{tabular}
\end{center}

规则：ICP > ISP 入栈，否则弹栈直到条件满足。

\begin{lstlisting}
void infixToPostfix() {
    Stack s; init(&s); char ch;
    while((ch=getchar())!='\n'){
        if(ch==' ')continue;
        if(ch>='0'&&ch<='9') printf("%c",ch);
        else if(ch=='(') push(&s,ch);
        else if(ch==')'){
            while(top(&s)!='(') printf("%c",pop(&s));
            pop(&s);  // 弹掉(
        } else { // + - * /
            while(!isEmpty(&s)&&isp(top(&s))>=icp(ch))
                printf("%c",pop(&s));
            push(&s,ch);
        }
    }
    while(!isEmpty(&s)) printf("%c",pop(&s));
}
\end{lstlisting}

\begin{bugbox}
忘了弹掉 \texttt{(}。忘了最后弹空栈。\\
只写 \texttt{push(ch)} 没有 \texttt{\&s}。
\end{bugbox}

\textbf{手算示例：} \texttt{(1+2)*3-4/(5+6)} \\
$\to$ \texttt{1 2 + 3 * 4 5 6 + / -}

\subsection*{栈应用三：函数调用}

每次函数调用创建\textbf{栈帧 (Stack Frame)}：返回地址、局部变量、旧帧指针。递归层数太深 $\to$ 栈溢出。

\textbf{尾递归消除：} 尾递归可转为迭代。例：\texttt{PrintList} 尾递归 → while 循环。

\section*{队列 (Queue) — FIFO}

\textbf{先进先出。} Enqueue(入队尾)、Dequeue(出队首)。

\subsection*{简单数组队列}

\begin{lstlisting}
typedef struct { int data[MAX]; int front,rear; } Queue;
void init(Queue* q) { q->front=0; q->rear=-1; }
void enq(Queue* q,int x) { q->data[++(q->rear)]=x; }
int deq(Queue* q) {
    if(q->front>q->rear)return -1;   // 判空
    return q->data[(q->front)++];
}
\end{lstlisting}

\textbf{缺点：} front 前面的空间永远浪费。

\subsection*{循环队列}

数组当环用。front=rear 为空，\texttt{(rear+1)\%MAX==front} 为满。\\
\textbf{浪费一个格子}区分空/满，最多存 MAX-1 个。

\begin{lstlisting}
void init(Queue* q) { q->front=0; q->rear=0; }
void enq(Queue* q, int x) {
    if((q->rear+1)%MAX==q->front) return;  // 先判满
    q->rear=(q->rear+1)%MAX;
    q->data[q->rear]=x;
}
int deq(Queue* q) {
    if(q->front==q->rear)return -1;        // 判空
    int x=q->data[q->front];
    q->front=(q->front+1)%MAX;
    return x;                               // return放最后
}
\end{lstlisting}

\begin{bugbox}
\texttt{init} 写成 \texttt{rear=-1}（循环队列应为 0）。\\
先移 rear 再判满 $\to$ 万一满的，rear 动了但数据没放。\\
dequeue 里 return 后还有 \texttt{front++} → 永远不会执行。\\
循环队列空用 \texttt{front>rear} → 应 \texttt{front==rear}。
\end{bugbox}

\subsection*{链表实现栈/队列}

栈：链表头插头删（头作栈顶）— $O(1)$。\\
队列：链表尾插头删（维护头尾指针）— $O(1)$。

\subsection*{层序遍历用队列}

二叉树的层序遍历：根入队 → 出队访问 → 左右孩子入队 → 重复。

\end{multicols}

\section*{线性结构复杂度总结}

\begin{center}
\begin{tabular}{|l|c|c|c|}
\hline
\textbf{操作} & \textbf{数组} & \textbf{单向链表} & \textbf{双向链表} \\
\hline
查找第k个 & $O(1)$ & $O(N)$ & $O(N)$ \\
头部插入 & $O(N)$ & $O(1)$ & $O(1)$ \\
尾部插入 & $O(1)$(有size) & $O(N)$(无尾针) & $O(1)$ \\
中间插入 & $O(N)$ & $O(1)$(已知前驱) & $O(1)$ \\
删除 & $O(N)$ & $O(1)$(已知前驱) & $O(1)$ \\
\hline
\end{tabular}
\end{center}

\end{document}
